Un tuyau d'orgue mesure \qty{\v{L}}{cm} de longueur. La température de l'air dans l'église est de \qty{\v{T}}{\degreeCelsius}. La vitesse du son dans l'air est donnée par $v = \qty{331}{m/s}\sqrt{T/\qty{273}{K}}$.

\begin{enumerate}[label=\alph*)]
    \item Déterminez la vitesse du son dans le tuyau.
    \item Si le tuyau est ouvert aux deux extrémités, déterminez les fréquences de ses trois premiers harmoniques.
    \item Si on bouche une extrémité du tuyau, déterminez les fréquences de ses trois premiers harmoniques. Pourquoi certains harmoniques disparaissent-ils ?
\end{enumerate}
